\documentclass[11pt,a4paper]{article}
\usepackage[T1]{fontenc}
\usepackage[utf8]{inputenc}
\usepackage[polish]{babel}
\usepackage[margin=2.5cm]{geometry}
\usepackage{amsmath,booktabs,hyperref}
\title{Regresja symboliczna w C++23\\Szkic raportu projektu}
\author{Studio projektowe 1}
\date{25 października 2026}

\begin{document}
\maketitle

\begin{abstract}
Projekt dotyczy implementacji regresji symbolicznej w C++23 oraz porównania
jej z PySR. Niniejszy dokument jest szkicem raportu: opisuje cel, podstawy
metody i zbiory danych. Obecny etap obejmuje przygotowanie środowiska,
generator danych, reprezentację wyrażeń i obliczanie ich wartości.
Wyszukiwanie ewolucyjne w C++ oraz końcowe porównanie nie są jeszcze wykonane.
\end{abstract}

\section{Cel i zakres}
Celem projektu jest wyszukiwanie interpretowalnych wzorów na podstawie danych
tabelarycznych. Planowane rozwiązanie obejmuje jedną populację ewolucyjną,
podstawową optymalizację stałych, konfigurowalne operatory i funkcję straty,
ograniczenia złożoności oraz proste szablony wyrażeń. Pytania badawcze dotyczą
odtwarzania zależności fizycznych w obecności szumu oraz kompromisu między
błędem testowym, złożonością i czasem działania względem PySR.

\section{Przegląd istniejących rozwiązań}
Regresja symboliczna poszukuje zarówno struktury modelu, jak i jego parametrów.
PySR wykorzystuje silnik SymbolicRegression.jl oraz wyszukiwanie ewolucyjne
łączone z upraszczaniem wyrażeń i optymalizacją stałych~\cite{cranmer}.
Oficjalne API udostępnia wybór operatorów, funkcji straty, limitów złożoności
i parametrów wyszukiwania~\cite{pysr}. Porównanie będzie dotyczyć całych
implementacji, a nie samej różnicy między językami C++ i Python.

\section{Metoda i bieżąca implementacja}
Wyrażenie jest drzewem z węzłami stałej, nazwanej zmiennej oraz operatora
unarnego lub binarnego. Obliczanie odbywa się rekurencyjnie z arytmetyką
\texttt{double}. API pozwala rejestrować własne operatory i funkcje straty;
dostępne funkcje straty to między innymi MSE i MAE. Niepoprawne dane,
nieznane zmienne lub operatory, dzielenie przez zero i wyniki NaN/Inf są
zgłaszane jako błędy. Testy obejmują wartości wzorcowe, zapis i odczyt CSV,
przypadki graniczne oraz powtarzalność generowania danych.

Na kolejnych etapach zostaną dodane inicjalizacja populacji, selekcja
turniejowa, mutacje, krzyżowanie i prosta optymalizacja stałych.

\section{Dane i metodologia oceny}
Wykorzystane zostaną trzy syntetyczne zależności fizyczne:
\begin{align*}
 F(x) &= 2.5x + 0.75, & x &\in [-2,2]\ \mathrm{m},\\
 h(t,v_0) &= v_0t - 4.905t^2, & t &\in [0,2]\ \mathrm{s},\quad v_0\in[12,20]\ \mathrm{m/s},\\
 x(t) &= 1.5\cos(2t+0.3), & t &\in [0,4]\ \mathrm{s}.
\end{align*}
Jednostkami wyjścia są odpowiednio N, m i m. Dane opisują idealizowane modele,
a nie rzeczywiste pomiary. Każdy problem ma 256 punktów treningowych,
128 walidacyjnych i 256 testowych, generowanych z niezależnymi ziarnami.
Wariant zaszumiony zachowuje te same wejścia i dodaje do celu szum
gaussowski o odchyleniu $0.01\,\mathrm{std}(y_{\mathrm{train,clean}})$.
Zapis CSV i manifest zawierają dane oraz parametry generowania.

Plan końcowy przewiduje 36 uruchomień: trzy zależności, dwa warianty danych,
trzy ziarna wyszukiwania i dwa systemy, z budżetem 90 sekund na wyszukiwanie.
Operatory, precyzja, limity złożoności, dane i zasoby CPU będą jednakowe.
Koszt startu i JIT zostanie oddzielony od wyszukiwania. Wzór zostanie wybrany
na danych walidacyjnych, a oceniony na testowych; raport obejmie błąd,
złożoność, czas i nieudane uruchomienia.

\section{Wyniki}
Na obecnym etapie rezultatem jest działająca reprezentacja i ewaluacja
zadanych wyrażeń oraz generator danych. Przykład sprężyny służy sprawdzeniu
obliczeń i API własnych funkcji; nie jest wynikiem automatycznego odkrycia wzoru.
Próbne uruchomienia PySR służą weryfikacji środowiska i eksportu wyniku,
nie zastępują końcowej serii porównawczej. Użyto PySR 2.7.2,
SymbolicRegression.jl 2.7.2, Julii 1.13.2 i Pythona 3.12.15, jednej populacji,
20 iteracji i ziarna 0. Wyboru wzoru dokonano według błędu walidacyjnego.
\begin{center}
\begin{tabular}{lrr}
\toprule
Próba & Złożoność & Testowe MSE \\
\midrule
Sprężyna, bez szumu & 5 & $2.82\cdot 10^{-29}$ \\
Ruch pionowy, bez szumu & 17 & $1.05\cdot 10^{-2}$ \\
Oscylator, z szumem & 18 & $1.09\cdot 10^{-4}$ \\
\bottomrule
\end{tabular}
\end{center}
Próba sprężyny zwróciła wzór bliski $2.5x+0.75$. Pojedyncze uruchomienia
nie stanowią oceny rozrzutu wyników ani dowodu przewagi nad implementacją C++.
Wyniki końcowej serii zostaną uzupełnione po implementacji wyszukiwania w C++.

\section{Dyskusja i ograniczenia}
Test poprawnego obliczania znanej formuły nie dowodzi skuteczności jej
wyszukiwania. Ograniczony zestaw danych i liczba powtórzeń nie pozwalają
na ogólne wnioski o przewadze jednego rozwiązania. Porównanie wymaga
zapisania wersji zależności, rzeczywistych danych i ustawień uruchomienia.

\section{Wnioski i dalsze prace}
Przygotowana baza będzie rozwijana o algorytm ewolucyjny i ograniczenia
wyrażeń. Wszystkie funkcje obowiązkowe są planowane do 30 listopada,
porównanie do 7 grudnia, a pełny pakiet i raport do 14 grudnia 2026 roku.

\begin{thebibliography}{9}
\bibitem{cranmer} M. Cranmer, \emph{Interpretable Machine Learning for Science
with PySR and SymbolicRegression.jl}, 2023.
\url{https://arxiv.org/abs/2305.01582}.
\bibitem{pysr} PySR, \emph{API Reference}.
\url{https://pysr.ai/v2.6.0/api}.
\end{thebibliography}
\end{document}
